\chapter{SOTTO LA PELLE}\label{sotto-la-pelle}

\hfill\break

Appresi molte cose sulla Perielio dai miei pazienti: soprattutto dagli
scienziati, a cui piaceva parlare, più che dagli amministratori, in
generale più taciturni; ma anche dalle famiglie del personale che
avevano iniziato ad abbandonare le loro cadenti compagnie di
assicurazione sanitaria a favore della clinica interna. All'improvviso
mi trovai a gestire un ambulatorio di medicina generale pienamente
funzionale, e quasi tutti i miei pazienti erano persone che avevano
osservato attentamente la realtà dello Spin e l'affrontavano con
coraggio e determinazione. «Il cinismo, qui, si ferma davanti al
cancello» mi disse un programmatore della missione. «Sappiamo che quello
che stiamo facendo è importante.» Era ammirevole. E anche contagioso.
Presto iniziai a considerarmi uno di loro e a sentire anche mio il
compito di estendere l'influenza umana al torrente impetuoso del tempo
extraterrestre.

Nei fine settimana risalivo la costa in auto fino a Kennedy per vedere
partire i razzi Atlas e Delta modernizzati che salivano rombando nel
cielo da una selva di piattaforme di lancio di nuova costruzione; e
talvolta, alla fine di quell'autunno o all'inizio dell'inverno, Jase
accantonava il lavoro e veniva con me. I missili erano semplici VRA,
dispositivi di ricognizione preprogrammati, finestre goffe sulle stelle.
I loro moduli di recupero sarebbero scesi lentamente (salvo un eventuale
fallimento della missione) nell'oceano Atlantico o sulle saline del
deserto occidentale, portando notizie dal mondo oltre il mondo.

Mi piaceva la grandiosità dei lanci. Ciò che affascinava Jase, ammise,
era la sconnessione relativistica che rappresentavano. Quelle piccole
casse di carico avrebbero potuto trascorrere settimane o addirittura
mesi oltre la barriera dello Spin, misurando la distanza dalla Luna che
si allontana o il volume del Sole in espansione, ma sarebbero cadute
sulla Terra (nel nostro sistema di riferimento) lo stesso pomeriggio,
bottiglie incantate riempite di un tempo superiore a quello che
avrebbero potuto contenere.

Quando quel vino fosse decantato, inevitabilmente le voci si sarebbero
rincorse per i corridoi della Perielio: le radiazioni gamma erano
aumentate, indice di qualche evento violento fra le stelle vicine;
c'erano nuove striature su Giove perché il Sole pompava più calore nella
sua atmosfera turbolenta; c'era un nuovo cratere enorme sulla Luna, che
non aveva più la faccia allineata con la Terra, ma che con una lenta
rotazione ci mostrava ora il suo lato oscuro.

Una mattina di dicembre Jase mi portò in un reparto tecnico dall'altra
parte del campus, dove era stato installato un modello dimostrativo a
grandezza naturale di un'astronave di carico marziana. Occupava una
piattaforma di alluminio in un angolo dell'enorme stanza a compartimenti
dove, intorno a noi, uomini e donne in tuta di Tyvek stavano assemblando
o allestendo per le prove altri prototipi. Il dispositivo era
incredibilmente piccolo, pensai. Era una scatola nera bitorzoluta delle
dimensioni di una cuccia con un ugello montato a un'estremità, sbiadita
dalle spietate luci del soffitto alto. Jase la ostentò con l'orgoglio di
un genitore.

«Fondamentalmente ha tre parti: il propulsore ionico con la massa di
reazione, i sistemi di navigazione di bordo e il carico utile. Quasi
tutta la massa è costituita dal motore. Non c'è nessun sistema di
comunicazione: non può inviare segnali alla Terra ma non serve che lo
faccia. I programmi di navigazione sono a ridondanza multipla ma
l'hardware non supera la grandezza di un telefono cellulare ed è
alimentato da pannelli solari.» I pannelli non erano installati ma
affissa alla parete c'era un'immagine artistica del veicolo assemblato
per intero, la cuccia trasformata in una libellula di Picasso.

«Non sembra abbastanza potente da riuscire a raggiungere Marte.»

«Non è la potenza il problema. I motori a ioni sono lenti ma saldi. Ed è
esattamente ciò che vogliamo: una tecnologia semplice, robusta e
duratura. La parte complicata è il sistema di navigazione, che
dev'essere intelligente e autonomo. Quando un oggetto supera la barriera
dello Spin, acquista quella che viene definita ``velocità temporale'',
una descrizione stupida ma che riesce a trasmettere l'idea. Il veicolo
di lancio viene accelerato e riscaldato, non rispetto a se stesso ma
rispetto a noi, e il differenziale è estremamente grande. Anche un
minuscolo cambiamento di velocità o traiettoria durante il lancio, una
piccola raffica di vento o un'alimentazione lenta del combustibile al
razzo vettore, rendono impossibile prevedere non come ma \emph{quando}
il veicolo comparirà nello spazio esterno.»

«Perché è importante?»

«È importante perché Marte e la Terra hanno entrambi un'orbita ellittica
e girano intorno al Sole a velocità diverse. Non esiste un modo
affidabile per calcolare in anticipo le relative posizioni dei pianeti
nel momento in cui il veicolo raggiunge l'orbita. In sostanza, la
macchina deve trovare Marte in un cielo affollato e tracciare la propria
traiettoria. Quindi abbiamo bisogno di un software intelligente e
flessibile e di un propulsore robusto e duraturo. Fortunatamente li
abbiamo entrambi. È una bella macchina, Tyler. Semplice esteticamente ma
grandiosa sotto la scocca. Prima o poi, lasciata ai suoi dispositivi e
salvo eventuali disastri, farà ciò per cui è stata progettata: entrerà
nell'orbita di Marte.»

«E poi?»

Jase sorrise. «È il cuore della questione. Guarda qui.» Smontò una serie
di bulloni finti dal prototipo e aprì un pannello sul davanti che rivelò
una camera schermata divisa in esagoni, a nido d'ape. In ogni esagono
era annidato un ovale nero senza punta. Un nido di uova color ebano.
Jason ne estrasse uno dalla sua sede. L'oggetto si poteva tenere in una
mano, tanto era piccolo.

«Sembra una freccetta incinta» ridacchiai.

«È solo un po' più sofisticato di una freccetta. Li spargeremo
nell'atmosfera marziana. Quando avranno raggiunto una determinata
altitudine, metteranno fuori le palette e nella discesa ruoteranno,
disperdendo calore e velocità. Dove spargerli - ai poli, all'equatore -
dipenderà dal carico di ciascun veicolo, il processo di base non
cambierà sia che cerchiamo depositi semiliquidi di acqua salata nel
sottosuolo o ghiaccio puro. Immaginali come aghi ipodermici che
iniettano la vita nel pianeta.»

Questa \emph{vita}, da quel che capivo, consisteva in microbi
geneticamente modificati, il cui materiale genetico era il frutto di una
combinazione tra batteri scoperti all'interno delle rocce delle valli
secche dell'Antartide, anaerobi in grado di sopravvivere nelle tubature
di scarico dei reattori nucleari e organismi unicellulari recuperati dal
fango ghiacciato sul fondo del mare di Barents. Questi organismi
avrebbero funzionato principalmente da ammendanti per il suolo, pensati
per svilupparsi mentre il Sole invecchiando riscaldava la superficie
marziana e rilasciava il vapore acqueo intrappolato e altri gas. In
seguito sarebbero arrivati ceppi ipermodificati di alghe azzurre,
semplici organismi autotrofi, e alla fine forme di vita più complesse in
grado di sfruttare l'habitat che i lanci iniziali avevano contribuito a
creare. Nella migliore delle ipotesi, Marte sarebbe rimasto comunque un
deserto; tutta l'acqua liberata avrebbe potuto creare al massimo qualche
lago instabile, salato e poco profondo\ldots{} ma forse sarebbe stato
sufficiente. Sufficiente a creare un luogo parzialmente abitabile oltre
la Terra avvolta nel suo sudario, dove gli esseri umani sarebbero andati
a vivere, un milione di secoli per ogni nostro anno. Dove i nostri
cugini marziani avrebbero avuto il tempo di risolvere enigmi cui noi
saremmo giunti solo brancolando.

Dove avremmo costruito, o permesso all'evoluzione di costruire per
nostro conto, una progenie di salvatori.

«È difficile credere che davvero possiamo fare tutto questo \ldots»

«\emph{Se} possiamo farlo. Non è per niente scontato.»

«Ma anche se lo potessimo fare, come soluzione è\ldots»

«È un atto di disperazione teleologica. Hai assolutamente ragione. Non
dirlo a voce alta, però. Comunque abbiamo una forza potente dalla nostra
parte.»

«Il tempo» supposi.

«No. Il tempo è una leva utile. Ma l'ingrediente attivo è la vita. La
vita in senso astratto, intendo: replicazione, evoluzione,
complessificazione. Il modo in cui la vita riempie crepe e fessure,
sopravvive facendo anche quello che non ti aspetti. Credo in questo
processo: è robusto, è saldo. Può salvarci? Non lo so. Ma la possibilità
è concreta.» Sorrise. «Se tu fossi il presidente di una commissione di
bilancio del Congresso sarei meno evasivo.»

Mi passò la freccetta. Era sorprendentemente leggera, non pesava più di
una palla da baseball della Major League. Provai a immaginarne centinaia
che piovevano da un cielo marziano senza nuvole, che fecondavano il
suolo sterile con il destino umano. Qualunque fosse il destino che ci
era rimasto.

\separatorefiori

E.D. Lawton visitò il complesso della Florida tre mesi dopo l'inizio del
nuovo anno, nello stesso periodo in cui si ripresentarono i sintomi di
Jason. Per qualche mese sembrava ci fosse stato un miglioramento.

Quando l'anno precedente Jase era venuto a trovarmi aveva descritto le
sue condizioni con riluttanza ma in maniera precisa. Debolezza e
intorpidimento transitori nelle braccia e nelle gambe. Vista annebbiata.
Vertigini episodiche. Incontinenza occasionale. Nessuno dei sintomi era
invalidante, ma erano diventati troppo frequenti per essere ignorati.

Poteva essere qualsiasi cosa, gli avevo detto, anche se doveva già
sapere quanto me che forse si trattava di un problema neurologico.

Provammo sollievo quando gli esami del sangue risultarono positivi alla
sclerosi multipla. La SM era una malattia curabile (o contenibile) da
quando dieci anni prima erano state introdotte le sclerostatine
chimiche. Una delle piccole ironie dello Spin era la sua coincidenza con
una serie di conquiste della medicina avvenute grazie alla ricerca
proteomica. La nostra generazione - quella di Jason e la mia - poteva
pure essere condannata, ma non sarebbe morta di SM, Parkinson, diabete,
cancro ai polmoni, arteriosclerosi o Alzheimer. L'ultima generazione del
mondo industrializzato sarebbe stata probabilmente la più sana.

Certo, non era proprio così semplice. Quasi il cinque per cento dei casi
diagnosticati di SM continuava a non rispondere al trattamento con le
sclerostatine né ad altre terapie. I clinici iniziavano a parlare di
casi di ``SMmultifarmaco-resistente'', forse si trattava persino di una
diversa malattia con la stessa sintomatologia.

Ma il trattamento iniziale di Jason aveva funzionato come previsto. Gli
avevo prescritto una dose giornaliera minima di Tremex e da allora si
era andato rimettendo completamente. Almeno fino alla settimana in cui
E.D. arrivò alla Perielio con tutta la grazia di una tempesta tropicale,
spargendo lungo i corridoi assistenti del Congresso e addetti stampa
come detriti spazzati dal vento.

E.D. era Washington, noi eravamo la Florida; lui era amministrazione,
noi scienza e ingegneria. Jase si manteneva in equilibrio un po'
precario tra i due fronti. Il suo lavoro consisteva essenzialmente nel
controllare che venissero rispettati i dettami del comitato direttivo,
ma si era spesso opposto alla burocrazia tanto che gli scienziati
avevano smesso di parlare di ``nepotismo'' e iniziato a offrirgli da
bere. Il guaio era, disse Jase, che E.D. non si accontentava di aver
dato il via al progetto Marte; voleva anche gestirlo nei minimi
dettagli, spesso per ragioni politiche, talvolta affidando contratti a
dubbi offerenti per comprarsi il sostegno del Congresso. Veniva deriso
dal personale, anche se quando era in città erano tutti abbastanza
contenti di stringergli la mano.

La visita ufficiale di quell'anno si concluse con un discorso al
personale e ai visitatori nell'auditorium del complesso. Entrammo tutti
in fila, diligenti come scolari ma più verosimilmente entusiasti, e non
appena l'uditorio si fu sistemato Jason si alzò per presentare suo
padre. Lo osservai salire sul palco e prendere il podio. Osservai il
modo in cui teneva la mano sinistra appesa vicino alla coscia, il modo
in cui si girò, ruotando goffamente sul tallone, quando strinse la mano
di suo padre.

Jase presentò suo padre brevemente ma in modo benevolo e poi si mescolò
di nuovo tra la folla di dignitari sul retro del palco. E.D. avanzò di
un passo. Aveva compiuto sessant'anni la settimana prima di Natale, ma
sarebbe potuto passare per un cinquantenne atletico, il ventre piatto
sotto un completo a tre pezzi e i capelli radi con un vivace taglio
militare a spazzola. Fece quello che sarebbe potuto essere un discorso
da campagna elettorale, lodò l'amministrazione Clayton per la sua
lungimiranza, il personale lì riunito per la sua dedizione alla
``visione Perielio'', suo figlio per la sua ``gestione ispirata'', gli
ingegneri e i tecnici per ``aver dato vita a un sogno e, se riusciremo,
per aver dato vita a un pianeta sterile e una nuova speranza a questo
mondo che continuiamo a chiamare casa''. Un'ovazione, un gesto di saluto
con la mano, un sorriso ferino, e non c'era già più, fatto sparire per
incanto dalla sua congrega di guardie del corpo.

Raggiunsi Jase un'ora più tardi nella mensa dei dirigenti, dove stava
seduto a un tavolino fingendo di leggere un estratto dell'
\emph{Astrophysics Review}.

Mi accomodai di fronte a lui. «Allora, come siamo messi?»

Sorrise debolmente. «Ti riferisci alla visita travolgente di mio padre?»

«Lo sai a cosa mi riferisco.»

Abbassò la voce. «Prendo i farmaci. Preciso come un orologio, mattina e
sera. Ma i sintomi sono tornati. Stamattina stavo male. Al braccio
sinistro, alla gamba sinistra, sentivo un formicolio. E sta peggiorando.
Molto più di prima. Quasi ogni ora. È come una scossa elettrica che mi
attraversa un lato del corpo.»

«Hai tempo di venire in infermeria?»

«Ho tempo, ma\ldots» Gli si inumidirono gli occhi. «Potrei non avere
modo. Non voglio allarmarti. Ma sono felice che tu sia arrivato. In
questo momento non sono sicuro di riuscire a camminare. Sono riuscito a
venire qui dopo il discorso di E.D. Ma sono certo che se provo ad
alzarmi cado. Non penso di riuscire a camminare. Ty\ldots{} \emph{non
	riesco a camminare}.»

«Vado a chiedere aiuto.»

Si drizzò sulla sedia. «No, non farai niente del genere. Se sarà
necessario, rimarrò seduto qui finché non vanno via tutti tranne la
guardia notturna.»

«È assurdo.»

«Oppure puoi aiutarmi ad alzarmi in maniera \emph{discreta. A} quanto
siamo, venti o trenta metri dall'infermeria? Se mi prendi per un braccio
e non ti mostri preoccupato, probabilmente possiamo arrivarci senza
attirare troppo l'attenzione.»

Alla fine accettai, non perché approvassi la farsa, ma solo perché
sembrava l'unico modo per portarlo nel mio ufficio. Gli presi il braccio
sinistro, lui puntellò la mano destra sul bordo del tavolo e si sollevò
a fatica. Riuscimmo ad attraversare la mensa senza andare a zigzag,
anche se Jason strascicava il piede sinistro in un modo difficile da
mascherare - fortunatamente nessuno ci guardò con molta attenzione. Una
volta raggiunto il corridoio, camminammo accostati al muro per rendere
meno evidente il suo trascinarsi. Quando un dirigente apparve in fondo
al corridoio, Jason sussurrò: «Fermati» e restammo lì come se stessimo
chiacchierando per caso, mentre Jason si reggeva a una teca e con la
mano destra stringeva lo scaffale in acciaio con una tale forza che le
nocche gli divennero bianche e sulla fronte comparvero evidenti perle di
sudore. Il dirigente ci passò accanto con un cenno muto.

Quando arrivammo all'entrata della clinica, ormai lo conducevo di peso.
Fortunatamente Molly Seagram non era in ufficio; una volta chiusa la
porta esterna, fummo soli. Lo aiutai a salire su un tavolo in uno degli
ambulatori, poi tornai al banco della reception e lasciai un appunto per
Molly per assicurarmi che non saremmo stati disturbati.

Quando tornai nell'ambulatorio, Jason stava piangendo. Non singhiozzava
ma le lacrime gli avevano rigato la faccia e si erano fermate sul mento.
«Cazzo, che cosa orribile.» Evitava di incrociare il mio sguardo. «Non
ce l'ho fatta, mi dispiace. Non ci sono riuscito.»

Aveva perso il controllo della vescica.

\separatorefiori

Lo aiutai a infilarsi un camice, gli sciacquai i vestiti bagnati nel
lavandino dell'ambulatorio e li misi ad asciugare accanto a una finestra
soleggiata nel deposito usato di rado, dietro gli armadi farmaceutici.
Quel giorno le cose da fare erano poche e fu quella la scusa che usai
per dare il pomeriggio libero a Molly.

Jason recuperò un po' della sua compostezza, anche se nel camice monouso
sembrava rimpicciolito. «Hai detto che si trattava di una malattia
curabile. Dimmi cos'è andato storto.»

«È guaribile, Jase. Per quasi tutti i pazienti, quasi sempre, ma ci sono
delle eccezioni.»

«E quindi io sono una di queste? Ho vinto la lotteria delle disgrazie?»

«È una ricaduta. È tipico delle malattie non curate, periodi di
invalidità si alternano a intervalli di remissione. Può anche darsi che
si tratti solo di una risposta lenta del tuo organismo. In alcuni casi
il farmaco, prima di essere pienamente efficace, deve raggiungere un
certo livello nel corpo e per un lungo periodo di tempo.»

«Sono passati sei mesi da quando me lo hai prescritto. E sto peggio, non
meglio.»

«Possiamo passare a un altro tipo di sclerostatine e vedere se la
situazione migliora, anche se chimicamente sono tutte molto simili.»

«Quindi cambiare prescrizione non servirà a niente.»

«Forse sì, forse no. Prima di scartare questa possibilità, proviamoci.»

«E se non funziona?»

«Allora smettiamo di provare a eliminare la malattia e iniziamo a
gestirla. La SM, anche se non curata, non è quasi mai una condanna a
morte. Un sacco di persone si rimettono completamente tra un attacco e
l'altro e riescono a condurre una vita relativamente normale.»

Anche se, non aggiunsi, di rado quei casi erano gravi o aggressivi come
sembrava essere quello di Jason. «Il consueto trattamento di riserva è
un cocktail di farmaci antinfiammatori, inibitori selettivi delle
proteine e stimolanti specifici del sistema nervoso centrale. Può essere
molto efficace nell'arrestare i sintomi e rallentare il decorso della
malattia.»

«Bene» disse Jason. «Fantastico. Scrivimi una ricetta.»

«Non è così semplice. Potresti andare incontro a effetti collaterali.»

«Per esempio?»

«Forse nessuno. Forse qualche disagio psichico: una lieve depressione o
episodi maniacali. Un po' di debolezza fisica generalizzata.»

«Ma sembrerò normale?»

«Presumibilmente sì.» ``Per ora,'' pensai, ``e magari per altri dieci o
quindici anni, forse di più.'' «Ma è una misura di controllo, non una
cura; un freno, non un arresto definitivo. Se avrai una vita abbastanza
lunga, la malattia tornerà.»

«Puoi darmi la certezza, comunque, che sarò a posto per una decina
d'anni?»

«La stessa certezza di tutte le cose di questa professione.»

«Dieci anni» mormorò pensieroso. «O un miliardo. Dipende da come la
guardi. Forse è sufficiente. Dovrebbe essere sufficiente, non credi?»

Non gli chiesi per cosa fosse sufficiente. Abbozzai: «Ma nel
frattempo\ldots»

«Non voglio un ``frattempo'', Tyler. Non posso permettermi di assentarmi
dal lavoro e non voglio che si sappia.»

«Non c'è nulla di cui vergognarsi.»

«Non me ne vergogno.» Indicò il camice con la mano destra. «Mi sento
umiliato a morte, ma non provo vergogna. Non è una questione
psicologica. Si tratta di quello che faccio qui alla Perielio. Di quello
che mi è \emph{concesso} fare. E.D. odia le malattie, Tyler. Odia
qualsiasi forma di debolezza. Ha odiato Carol dal giorno in cui il bere
è diventato un problema.»

«Non credi che capirebbe?»

«Voglio bene a mio padre, ma conosco i suoi difetti. No, non capirebbe.
L'influenza che ho qui passa per lui. E al momento è un po' precaria.
Abbiamo avuto dei dissapori. Se per lui diventassi un peso, prima della
fine della settimana mi farebbe confinare in qualche costosa casa di
cura in Svizzera o a Bali, e si racconterebbe che l'ha fatto per il mio
bene. E quel che è peggio, ci crederebbe pure.»

«Quello che scegli di rendere pubblico non mi riguarda. Ma devi vedere
un neurologo, non un medico generico.»

«No.»

«In coscienza, Jase, non posso continuare a curarti se non parli con uno
specialista. È stato già abbastanza azzardato prescriverti il Tremex
senza consultare prima un esperto del cervello.»

«Hai la risonanza magnetica e gli esami del sangue, giusto? Cos'altro ti
serve?»

«L'ideale sarebbe un laboratorio ospedaliero completamente attrezzato e
una laurea in neurologia.»

«Cazzate. L'hai detto tu stesso, oggigiorno la SM non è un grosso
problema.»

«A meno che non risponda alla terapia.»

«Non posso\ldots» Aveva voglia di discutere, ma era visibilmente
stanchissimo. L'affaticamento poteva essere un altro sintomo della sua
ricaduta, comunque; aveva lavorato troppo nelle settimane precedenti la
visita di E.D. «Facciamo un patto. Vedrò uno specialista se riesci a
organizzare la cosa con discrezione e a non tenerne traccia in alcuna
cartella alla Perielio. Ma devo essere operativo. Devo essere operativo
\emph{domani}. Operativo nel senso che devo poter camminare senza aiuto
e non devo pisciarmi addosso. Il cocktail di farmaci di cui mi hai
parlato fa effetto subito?»

«Di solito sì. Ma senza un controllo neurologico\ldots»

«Tyler, devo dirtelo, apprezzo quello che hai fatto per me, ma se ne ho
bisogno posso comprarmi un dottore più collaborativo. Dammi subito
questa terapia e vedrò uno specialista, farò tutto quello che pensi sia
giusto fare. Ma se credi che mi presenti al lavoro su una sedia a
rotelle e con un catetere nell'uccello, ti sbagli di grosso.»

«Anche se ti faccio la prescrizione, non è che domani starai già meglio,
Jase. Ci vogliono un paio di giorni.»

«A un paio di giorni potrei pure rinunciare.» Sembrò pensarci su. Quindi
cedette: «Va bene, voglio i farmaci e voglio che mi tiri fuori di qui
senza dare nell'occhio. Se ci riesci, sono nelle tue mani. E niente
discussioni.»

«I medici non contrattano, Jase.»

«Prendere o lasciare, Ippocrate.»

\separatorefiori

Non gli diedi subito tutto il cocktail di farmaci - la nostra farmacia
non ne era provvista - ma cominciai con uno stimolante dell'SNC che gli
avrebbe restituito almeno il controllo della vescica e la capacità di
camminare senza assistenza per i giorni successivi. Il rovescio della
medaglia era uno stato mentale irritabile, gelido, simile - almeno così
mi era stato detto - a quello di chi usciva da un fine settimana dedito
all'abuso di cocaina. Gli alzò la pressione sanguigna e gli fece
comparire delle borse scure sotto gli occhi.

Aspettammo che quasi tutto il personale andasse a casa e che nel
complesso ci fossero solo quelli del turno di notte. Quando superò la
reception per arrivare al parcheggio Jase camminava rigido ma in maniera
credibile, salutò amichevolmente un paio di colleghi che uscivano in
ritardo e affondò nel sedile del passeggero della mia auto. Lo portai a
casa sua.

Era venuto a trovarmi diverse volte nella mia casetta in affitto, ma io
non ero mai stato a casa sua. Mi aspettavo qualcosa che riflettesse il
suo status alla Perielio. In effetti, l'appartamento in cui dormiva -
era evidente che ci facesse poco altro - era una modesta unità
condominiale con un pezzetto di vista sull'oceano. Lo aveva arredato con
un divano, un televisore, una scrivania, un paio di librerie e una
connessione a banda larga. Le pareti erano spoglie tranne lo spazio
sopra la scrivania, dove aveva attaccato con il nastro adesivo un
diagramma disegnato a mano che raffigurava la storia lineare del sistema
solare dalla nascita del Sole al suo collasso finale in una nana bianca
ardente, con la storia umana che divergeva dalla linea in un punto
contrassegnato come {LO} S{PIN}. Le librerie erano colme di riviste e
testi accademici, e decorate con tre sole fotografie incorniciate: E.D.
Lawton, Carol Lawton e un'immagine dimessa di Diane che doveva essere
stata scattata anni prima.

Jase si stese sul divano. Sembrava il ritratto di un paradosso, con il
corpo a riposo e gli occhi che brillavano ipervigili per via dei
farmaci. Andai nella piccola cucina adiacente e preparai delle uova
strapazzate (nessuno dei due aveva mangiato dopo la colazione) mentre
Jason parlava. E parlava ancora. E continuava a parlare. A un certo
punto disse: «Lo so che sto parlando troppo, ne sono consapevole, ma non
riesco nemmeno a pensare di dormire: è una cosa che passa questa?»

«Se ti faremo assumere il cocktail di farmaci per molto tempo, allora
sì, l'effetto stimolante svanirà.» Gli portai un piatto sul divano.

«È molto veloce. Sembra una di quelle pillole che la gente prende quando
deve prepararsi in fretta per gli esami finali. Fisicamente, però,
rilassa. Mi sento come un'insegna al neon su un edificio vuoto. Tutto
illuminato ma senza alcun valore. Le uova, le uova sono molto buone.
Grazie.» Scostò il piatto. Ne aveva mangiato forse un cucchiaio.

Mi sedetti alla sua scrivania e diedi un'occhiata al diagramma dello
Spin sulla parete opposta. Mi chiesi come fosse convivere con questa
rappresentazione cruda delle origini dell'umanità e del suo destino, la
specie umana ridotta a un evento finito nella vita di una stella
ordinaria. Lo aveva disegnato con un pennarello su un comune rotolo di
carta da imballaggio marrone.

Jason seguì il mio sguardo. «Evidentemente vogliono che facciamo
\emph{qualcosa}\ldots»

«Chi?»

«Gli Ipotetici. Se dobbiamo chiamarli così. E suppongo di sì. Tutti li
chiamano così. Si aspettano qualcosa da noi. Non so cosa. Un dono, un
segnale, un sacrificio accettabile.»

«Come fai a saperlo?»

«Non è certo un'osservazione originale la mia. Perché la barriera dello
Spin è permeabile agli artefatti umani, come i satelliti, ma non alle
meteore e neanche alle particelle Brownlee? Perché evidentemente
\emph{non} è una barriera; non è mai stata la parola giusta per
definirla.» Sotto l'effetto dello stimolante Jase sembrava
particolarmente affezionato alla parola \emph{evidentemente}.
«Evidentemente è un filtro selettivo. Sappiamo che filtra l'energia che
arriva sulla superficie terrestre. Quindi gli Ipotetici vogliono
mantenerci in vita e intatti, o perlomeno l'ecologia terrestre, ma
allora perché concederci di accedere allo spazio? Anche dopo aver
tentato di distruggere gli unici due artefatti legati allo Spin che
siano mai stati trovati? Cosa stanno \emph{aspettando}, Ty? Qual è il
premio?»

«Forse non è un premio. Forse è un riscatto. ``Pagate e vi lasceremo in
pace.''»

Scosse la testa. «È troppo tardi per lasciarci in pace. Abbiamo bisogno
di loro. E non possiamo ancora scartare la possibilità che siano
benevoli o almeno amichevoli. Poniamo che non fossero intervenuti nel
momento in cui l'hanno fatto. Cosa ci sarebbe toccato? Tanti dicevano
che la nostra società non era più sostenibile e che eravamo giunti al
capolinea addirittura come specie. Riscaldamento globale,
sovrappopolazione, morte dei mari, perdita di terre coltivabili,
proliferazione di malattie, minacce di guerre nucleari o
biologiche\ldots»

«Magari ci saremmo distrutti con le nostre stesse mani ma almeno sarebbe
stata colpa nostra.»

«Lo sarebbe stato davvero? Colpa \emph{di chi}, esattamente? Tua? Mia?
No, sarebbe stato il risultato di diversi miliardi di esseri umani che
hanno fatto scelte relativamente innocue: avere figli, andare a lavoro
in auto, cercare di mantenere quel lavoro, risolvere prima di tutto i
problemi dell'immediato futuro. Quando arrivi al punto in cui anche le
azioni più banali sono punibili con la morte della specie, allora
evidentemente, evidentemente, sei in una fase critica, un tipo diverso
di punto di non ritorno.»

«È meglio essere consumati dal Sole?»

«Non è ancora successo. E la nostra non è la prima stella che si spegne.
La galassia è disseminata di nane bianche che un tempo avrebbero potuto
ospitare pianeti abitabili. Ti chiedi mai cosa sia accaduto
\emph{loro}?»

«Raramente» risposi.

Mi avvicinai alla libreria, alle foto di famiglia, camminando sul nudo
pavimento in parquet. C'era E.D. che sorrideva alla macchina
fotografica: un uomo i cui sorrisi non erano mai del tutto convincenti.
La somiglianza fisica con Jason era notevole. (``Evidente'', avrebbe
detto forse Jase.) Macchina simile, spirito differente.

«Come potrebbe sopravvivere la vita a una catastrofe stellare? Ma
dipende da cosa si intende per vita, evidentemente. Parliamo di vita
organica o di qualsiasi tipo di ciclo di retroazione autocatalitico
generalizzato? Gli Ipotetici sono organici? Una domanda interessante di
per sé\ldots»

«Dovresti proprio cercare di dormire un po'.»

Era mezzanotte passata e aveva iniziato a usare parole che non capivo.
Presi la foto di Carol. Qui la somiglianza con Jason era più
impercettibile. Il fotografo l'aveva immortalata in una giornata
positiva: aveva gli occhi aperti, non bloccati a mezz'asta, e anche se
sorrideva di malavoglia, le labbra sottili sollevate in modo appena
percettibile, il sorriso non era del tutto finto.

«Forse stanno minando il Sole» disse Jason, continuando a parlare degli
Ipotetici. «Disponiamo di alcuni dati suggestivi sui brillamenti solari.
Evidentemente, quello che hanno fatto alla Terra richiede enormi
quantità di energia utilizzabile. È come raffreddare una massa di
dimensioni planetarie a una temperatura prossima allo zero assoluto. E
l'alimentatore dov'è? Molto probabilmente è il Sole. E dopo lo Spin
abbiamo osservato una riduzione notevole dei grandi brillamenti solari.
Qualcosa, una forza o un'azione, forse sfrutta le particelle ad alta
energia prima che salgano verso l'eliosfera. Minano il Sole, Tyler! È un
atto di superbia tecnologica sorprendente quasi quanto lo stesso Spin.»

Presi la foto incorniciata di Diane. Era precedente al suo matrimonio
con Simon Townsend. Il fotografo aveva catturato una certa inquietudine
particolare, come se avesse appena socchiuso gli occhi per un pensiero
sconcertante. Era bella naturalmente ma non proprio a suo agio,
armoniosa ma allo stesso tempo un po' fuori equilibrio.

Avevo tantissimi ricordi di lei. Ma quei ricordi erano ormai vecchi di
anni, svanivano nel passato con una velocità quasi simile a quella dello
Spin. Jason vide che reggevo la cornice e rimase in silenzio per qualche
istante benedetto. Poi sbottò: «Credimi, Tyler, questa tua fissazione è
indegna.»

«Non è affatto una fissazione, Jase.»

«Perché? Perché l'hai \emph{superata} o perché ne hai \emph{paura}?
Potrei fare anche a lei la stessa domanda. Se chiamasse qualche volta.
Simon la tiene al guinzaglio. Sospetto che rimpianga i giorni ai tempi
dell'NR, quando il movimento era pieno di hippy unitariani ed
evangelici. Il prezzo della religiosità adesso è più salato.» Aggiunse:
«Di tanto in tanto parla con Carol.»

«Almeno è felice?»

«Diane si trova tra fanatici. Potrebbe esserlo anche lei. La felicità
non è contemplata.»

«Pensi sia in pericolo?»

Si strinse nelle spalle. «Penso stia vivendo la vita che ha scelto per
se stessa. Avrebbe potuto fare altre scelte. Per esempio, avrebbe potuto
sposare \emph{te}, Ty, se non fosse stato per quella sua ridicola
fantasia\ldots»

«Che fantasia?»

«Che E.D. è tuo padre. Che è la tua sorella biologica.»

Indietreggiai dalla libreria troppo in fretta e quindi urtai le foto che
caddero a terra.

«È ridicolo.»

«Chiaramente ridicolo. Credo che abbia abbandonato del tutto l'idea solo
quando è andata al college.»

«Come ha potuto anche solo pensare\ldots»

«Era una fantasia, non una teoria. Pensaci. Non c'è mai stato molto
affetto tra Diane e E.D. Si sentiva ignorata da lui. E in un certo
senso, aveva ragione. E.D. non ha mai voluto una figlia, voleva un
erede, un erede maschio. Aveva grandi aspettative e il caso ha voluto
che io ne fossi all'altezza. Diane per lui era una distrazione. Si
aspettava che se ne occupasse Carol, ma Carol\ldots» scrollò le spalle.
«Carol non era all'altezza del compito.»

«E così si è inventata questa\ldots{} storia?»

«Lei la considerava una deduzione. Spiegava perché E.D. vi facesse
vivere nella proprietà. Spiegava la costante infelicità di Carol. E, in
fondo, la faceva sentire bene con se stessa. Tua madre nei suoi
confronti era più gentile e premurosa di quanto lo fosse mai stata
Carol. Le piaceva l'idea di essere imparentata con la famiglia Dupree.»

Guardai Jason. Aveva la faccia pallida, le pupille dilatate, lo sguardo
distante e rivolto alla finestra. Ricordai a me stesso che era un mio
paziente, che manifestava una reazione psicologica prevedibile a un
farmaco potente; che quello era lo stesso uomo che solo qualche ora
prima aveva pianto per la propria incontinenza. Dissi: «Ora devo proprio
andarmene, Jason.»

«Perché? Ti sciocca così tanto tutto questo? Pensavi che crescere
dovesse essere indolore?» Poi, all'improvviso, prima che potessi
rispondere, girò la testa e per la prima volta quella sera incrociò il
mio sguardo. «Accidenti, forse mi sto comportando in modo pessimo.»

«Le medicine\ldots»

«Mi sto comportando in modo orrendo. Mi dispiace, Tyler.»

«Ti sentirai meglio dopo una bella dormita. Ma per un paio di giorni non
devi tornare alla Perielio.»

«No, non ci vado. Fai un salto domani?»

«Sì.»

«Grazie.»

Me ne andai senza rispondere.
